\documentclass[11pt]{amsart}
\usepackage[utf8]{inputenc}
\usepackage[T1]{fontenc}
\usepackage{amssymb}
\usepackage{booktabs}
\usepackage[margin=1.2in]{geometry}
\usepackage[hidelinks]{hyperref}

\newtheorem{theorem}{Theorem}[section]
\newtheorem{lemma}[theorem]{Lemma}
\newtheorem{proposition}[theorem]{Proposition}
\newtheorem{corollary}[theorem]{Corollary}
\newtheorem*{theoremA}{Theorem A}
\newtheorem*{theoremB}{Theorem B}
\newtheorem{mainthm}{Theorem}
\renewcommand{\themainthm}{\arabic{mainthm}}
\theoremstyle{definition}
\newtheorem{definition}[theorem]{Definition}
\newtheorem{question}[theorem]{Question}
\theoremstyle{remark}
\newtheorem{remark}[theorem]{Remark}

\newcommand{\F}{\mathbb{F}_2}
\newcommand{\GL}{\mathrm{GL}}
\newcommand{\SL}{\mathrm{SL}}
\newcommand{\rk}{\operatorname{rank}}
\newcommand{\Span}{\operatorname{span}}
\newcommand{\cn}{\operatorname{cn}}

\title{Kourovka Problem 16.21 over the field with two elements}
\author{AdityaVG13}
\date{}

\begin{document}

\begin{abstract}
L.~Vaserstein asked (Kourovka Notebook, Problem 16.21) whether, for $n>2$ and a non-central
matrix $\alpha\in\SL_n(F)$, every non-central matrix in $\SL_n(F)$ is a product of $n$ matrices,
each similar to $\alpha$. We answer the question affirmatively for $F=\F$, and show that $n$ may
be replaced by any $k\ge n$. The proof assumes two published theorems, of Kn\"uppel and Nielsen
and of B\"unger and Nielsen. It is an induction on $n$ that sorts matrices by one property:
whether they fix a vector together with a complementary hyperplane. Matrices that do reduce to
dimensions $n-1$ and $n-2$ through an explicit construction. Matrices that do not have minimum
rank at least $3$, and for them the two cited theorems give a product of at most
$2\lceil(n-1)/3\rceil+1$ conjugates, about $2n/3$. The number $n$ is needed only because of
matrices close to the identity, such as transvections. We give the exact number of conjugates
needed for every conjugacy class of $\GL_3(\F)$ and $\GL_4(\F)$. Everything except the two cited
theorems is formalized in Lean~4 with Mathlib. The work was produced almost entirely by an AI agent
under the author's direction, as a test of a new computational mathematics project, and we
describe how.
\end{abstract}

\maketitle

\section{Introduction}

\subsection{The question}
Problem 16.21 of the Kourovka Notebook \cite{Kourovka}, posed by L.~Vaserstein, reads:

\begin{quote}
Given a non-central matrix $\alpha\in\SL_n(F)$ over a field $F$ for $n>2$, is it true that every
non-central matrix in $\SL_n(F)$ is a product of $n$ matrices, each similar to $\alpha$?
\end{quote}

The 21st edition of the Notebook (2026) lists the problem as open.

The question belongs to a long line of work on products of conjugacy classes in groups; see the
volume of Arad and Herzog \cite{AH} for the classical theory. For a conjugacy class $C$ of a group
$G$, write $C^k$ for the set of products of $k$ elements of $C$. A typical question asks how large
$k$ must be for $C^k$ to contain all of $G$, or all elements outside the centre. For linear groups
this is a question about matrices: given one matrix $\alpha$, how many matrices similar to $\alpha$
does it take to write an arbitrary matrix as their product?

The number $n$ in Vaserstein's question is natural. If $\alpha$ is a transvection, every matrix
$c$ similar to $\alpha$ has $\rk(c-1)=1$. From $xy-1=x(y-1)+(x-1)$ we get
$\rk(xy-1)\le\rk(x-1)+\rk(y-1)$, so a product of $k$ transvections $g$ has $\rk(g-1)\le k$. A
matrix $g$ with $g-1$ invertible therefore needs at least $n$ transvections. So $n$ cannot be
lowered in general, and the question asks whether it is always enough.

\subsection{What was known}
For fields with at least four elements, Lev \cite{Lev} showed that a product of two prescribed
cyclic similarity classes of $\GL_n(F)$ contains every nonscalar matrix of the right determinant,
under a condition on eigenvalues. Kn\"uppel \cite{Knuppel} studied the least $k$ such that every
element of $\SL(V)$ is a product of at most $k$ elements of a given class. Kn\"uppel and Nielsen
\cite{KN} determined the \emph{extended covering number} of $\SL_n(K)$, the least $k$ such that
$1\in C_1\cdots C_k$ for all non-central classes $C_1,\dots,C_k$; it equals $n+1$ ``provided $n$
and the field are not too small''. The field $\F$ is the smallest field, and parts of their
argument do not apply to it.

We do not know whether these results settle Problem 16.21 for any field, and we make no claim
about fields other than $\F$.

\subsection{Results}
Over $\F$ every invertible matrix has determinant $1$, so $\SL_n(\F)=\GL_n(\F)$, the centre is
trivial, and a matrix similar to $\alpha$ is the same as a conjugate of $\alpha$ in $\GL_n(\F)$.
The question becomes: is every nonidentity element of $\GL_n(\F)$ a product of $n$ conjugates of
any given nonidentity element?

\begin{mainthm}\label{thm:main}
Assume Theorems A and B below. Let $n\ge3$, let $\alpha,g\in\GL_n(\F)$ with $\alpha\ne1$ and
$g\ne1$, and let $k\ge n$. Then $g$ is a product of $k$ conjugates of $\alpha$. In particular,
Problem 16.21 has a positive answer for $F=\F$.
\end{mainthm}

The proof shows more for matrices that are far from the identity. Say that $h\in\GL(V)$ has a
\emph{J1 block} if it fixes a vector $v$ and a linear form $\varphi$ with $\varphi(v)=1$; then $h$
is the direct sum of the identity on $\langle v\rangle$ and a map on the hyperplane $\ker\varphi$.
In terms of the Jordan form, $h$ has a Jordan block of size $1$ for the eigenvalue $1$.

\begin{mainthm}\label{thm:far}
Assume Theorems A and B. Let $n\ge5$ and let $\alpha\in\GL_n(\F)$, $\alpha\ne1$, have no J1 block.
Then every nonidentity element of $\GL_n(\F)$ is a product of $k$ conjugates of $\alpha$ for every
$k\ge 2\lceil(n-1)/3\rceil+1$.
\end{mainthm}

So for most matrices about $2n/3$ conjugates suffice, and $n$ is forced only by matrices that
split off an identity block. For $n=3$ and $n=4$ we computed the exact numbers
(Section~\ref{sec:small}, Table~\ref{tab:cn}). Most classes need only $2$ conjugates, and the
classes that need $n$ are the transvections and, for $n=4$, one class of elements of order~$3$.

\subsection{The cited theorems}
A matrix $A\in K^{n\times n}$ is \emph{cyclic} if there is a vector $v$ such that
$v,Av,\dots,A^{n-1}v$ is a basis of $K^n$. The \emph{minimum rank} of $A$ is
$\delta(A)=\min_{\lambda\in\overline K}\rk(A-\lambda)$, where $\overline K$ is an algebraic
closure of $K$. Kn\"uppel and Nielsen define $\delta(A)$ as $n$ minus the number of invariant
divisors of $A$ different from $1$, and note in the proof of their Lemma~5.1 that the two
definitions agree. Both notions depend only on the similarity class.

\begin{theoremA}[{Kn\"uppel and Nielsen \cite[Corollary 6.9]{KN}, case $K=\F$}]
Let $n\ge3$, $n\ne4$, and let $C_1,\dots,C_m$ be non-central conjugacy classes of $\SL_n(\F)$ with
$\delta(C_1)+\dots+\delta(C_m)\ge n-1$. Then $C_1C_2\cdots C_m$ contains a cyclic matrix.
\end{theoremA}

\begin{theoremB}[{B\"unger and Nielsen \cite[Theorem 1.2]{BN}, case $K=\F$}]
Let $n\ge3$, let $M\in\GL_n(\F)$ be a matrix that is not a homothety, let $X$ be a cyclic
similarity class of $\GL_n(\F)$, and let $q\in\F[x]$ be monic of degree $n-1$ with $q(0)\ne0$.
Then there is a similarity class $D$ of $\GL_n(\F)$ whose minimal polynomial is divisible by $q$
and such that $M\in DX$.
\end{theoremB}

We do not prove these two theorems. Every other step is elementary, and every other step is
checked in Lean~4 \cite{Lean4} with Mathlib \cite{Mathlib}: the formal main theorem takes
Theorems A and B as hypotheses and depends on no other assumptions (Section~\ref{sec:lean}).

\subsection{How the result was found}
This work was produced almost entirely by an AI agent (Droid, by Factory), working under the
author's direction. The author is developing a new computational mathematics project, a
programming language for exact mathematical computation in which one writes definitions as on
paper and states expected facts as checks that the program evaluates, and this problem was taken
up to test that project on open problems. The agent carried out the computations, found the
argument, searched the literature, wrote the Lean formalization and drafted this text. The author
chose the problem, set the scope (in particular, to cite Theorems A and B rather than formalize
them) and takes responsibility for the paper. No specialist has refereed the argument; the main
evidence for its correctness is the Lean check of every step outside Theorems A and B.

The path to the proof went as follows.

\emph{Small cases.} An exhaustive computation in $\GL_3(\F)$ and $\GL_4(\F)$ confirmed the answer
for $n=3,4$ and produced the data of Table~\ref{tab:cn}. Random searches for explicit products
found witnesses for every pair of classes for $n=5$ and $n=6$. This suggested that the answer is
yes and that most classes have a lot of room to spare.

\emph{Matrices with a J1 block.} If $\alpha=1\oplus\alpha'$, induction handles every target $g$
that also has a J1 block. The data suggested that, for every other target, one conjugate $c$ of
$\alpha$ can be chosen so that $gc^{-1}$ has a J1 block; this became Lemma~\ref{lem:P}. Testing a
construction of $c$ on every pair of classes for $n=4,5,6$ succeeded outside one situation:
$\alpha$ an involution and $g$ satisfying $g^2+g+1=0$. In $\GL_4(\F)$, none of the $315$ involutions works for
such $g$. Remark~\ref{rem:sharpP} explains why, and Lemma~\ref{lem:exception} handles this case
directly.

\emph{Matrices without a J1 block.} Here the reduction does not apply, and several attempts at a
direct argument stalled. A literature search found Theorems A and B, which together turn
``$\alpha$ has large minimum rank'' into ``every $g$ is a short product of conjugates''. What
remained was to show that matrices without a J1 block have minimum rank at least $3$
(Lemma~\ref{lem:minrank}).

\emph{Formalization.} The first complete argument split the classes by minimum rank and used a
third published result for transvections. Writing the proof in Lean led to the simpler split used
here, by whether $\alpha$ has a J1 block, and the third result was no longer needed.

\subsection{Organization}
Section~\ref{sec:overview} explains the idea of the proof. Section~\ref{sec:prelim} collects
preliminaries, and Section~\ref{sec:small} treats $n=3,4$ and gives the exact data.
Sections~\ref{sec:J1} and~\ref{sec:noJ1} treat matrices with and without a J1 block, and
Section~\ref{sec:proof} completes the proofs. Section~\ref{sec:lean} describes the formalization
and Section~\ref{sec:discussion} lists open questions.

\section{The idea of the proof}\label{sec:overview}

Fix $\alpha\ne1$ and write $C$ for its class. Two simple facts organize everything.

First, adding factors never hurts: once every nonidentity element lies in $C^m$, the same holds
for $C^k$ with $k\ge m$ (Lemma~\ref{lem:more}). So it is enough to reach some $m\le n$.

Second, matrices come in two kinds. A matrix with a J1 block is ``the identity plus something
smaller''. A matrix without one is spread out: no eigenvalue can have a large eigenspace.

\emph{If $\alpha$ has a J1 block}, write $\alpha=1\oplus\alpha'$. A target of the same shape,
$g=1\oplus g'$, is handled in dimension $n-1$ by induction. For any other target $g$, we choose one
conjugate $c$ of $\alpha$ so that $gc^{-1}$ has a J1 block. Then $gc^{-1}$ is a product of $n-1$
conjugates by the previous case, and $g=(gc^{-1})c$ is a product of $n$. Finding $c$ is a matter of
lining up two pairs of vectors and linear forms, one adapted to $g$ and one to $\alpha$. This works
unless $\alpha$ is an involution and $g$ is a sum of rotations of order $3$, and that case is done
by hand, using dimension $n-2$.

\emph{If $\alpha$ has no J1 block}, its minimum rank is at least $3$. Theorem~A then says that a
product of about $(n-1)/3$ conjugates of $\alpha$ is a cyclic matrix $D$. Theorem~B says that
every non-homothety factors as (a matrix of a special class $N$) times (a conjugate of $D$). The
class $N$ is the same whatever matrix is factored, and applying Theorem~B to $\alpha$ itself shows
that $N\subseteq C\cdot C^{k}$. So every $g$ lies in $N\cdot C^k\subseteq C^{2k+1}$, and
$2k+1\le n$.

The base cases $n=3$ and $n=4$ are done by computer, because the induction needs dimensions $n-1$
and $n-2$ and Theorem~A excludes $n=4$.

\section{Preliminaries}\label{sec:prelim}

Throughout, $V$ is a vector space of dimension $n$ over $\F$, $V^*$ is its dual and
$G=\GL(V)$. For $a\in G$ let $C_a$ be its conjugacy class and let $C_a^k$ be the set of products
of $k$ elements of $C_a$. The set $C_a^k$ is closed under conjugation and depends only on $C_a$.
For $n\ge2$ the centre of $G$ is trivial, so every $a\ne1$ is non-central.

For $n\ge 1$ let $P(n)$ be the statement: for every space $V$ of dimension $n$ and all
$a,g\in\GL(V)\setminus\{1\}$, we have $g\in C_a^k$ for every $k\ge n$.

\begin{lemma}\label{lem:more}
Let $a\in G$ be non-central and suppose $G\setminus\{1\}\subseteq C_a^m$. Then
$G\setminus\{1\}\subseteq C_a^k$ for every $k\ge m$.
\end{lemma}

\begin{proof}
Let $k=m+j$ with $j\ge1$ and $g\ne1$. The set $C_a^j$ contains $a^j$ and $(xax^{-1})a^{j-1}$, where
$xax^{-1}\ne a$, and these differ. So some $u\in C_a^j$ satisfies $u\ne g$. Then $u^{-1}g\ne1$
lies in $C_a^m$, and $g=u\cdot u^{-1}g\in C_a^{j+m}$.
\end{proof}

Over $\F$, vectors $x_1,\dots,x_m$ are linearly independent if and only if no nonempty subset of
them sums to $0$. A vector space over any field is not the union of two proper subspaces. If
$x_1,\dots,x_m\in V$ are linearly independent, there is $\varphi\in V^*$ with any prescribed values
$\varphi(x_i)$.

\begin{definition}
An element $h\in G$ has a \emph{J1 block} if there are $v\in V$ and $\varphi\in V^*$ with
$\varphi(v)=1$, $hv=v$ and $\varphi\circ h=\varphi$.
\end{definition}

If $h$ has a J1 block then $V=\langle v\rangle\oplus\ker\varphi$, both summands are
$h$-invariant, and $h=1\oplus h'$ with $h'\in\GL(\ker\varphi)$. Having a J1 block is invariant
under conjugation. With $T=h-1$, a J1 block exists exactly when $\ker T\not\subseteq\operatorname{im}T$,
that is, when $n-2\rk T+\rk T^2>0$.

A \emph{frame} of size $m$ is a pair of families $e_1,\dots,e_m\in V$ and
$f_1,\dots,f_m\in V^*$ with $f_i(e_j)=\delta_{ij}$. Then
$V=\Span(e_1,\dots,e_m)\oplus\bigcap_i\ker f_i$.

\begin{lemma}\label{lem:frame}
Given two frames $(e,f)$ and $(e',f')$ of the same size $m$, there is $t\in G$ with $te'_j=e_j$
and $f_i\circ t=f'_i$ for all $i,j$.
\end{lemma}

\begin{proof}
Both kernels $\bigcap_i\ker f_i$ and $\bigcap_i\ker f'_i$ have dimension $n-m$. Let $t$ send
$e'_j$ to $e_j$ and map the second kernel isomorphically onto the first. For
$x=\sum_j f'_j(x)e'_j+y$ with $y$ in the second kernel, $tx=\sum_j f'_j(x)e_j+ty$, so
$f_i(tx)=f'_i(x)$.
\end{proof}

We also use the standard fact that if $a=a_1\oplus a_2$ along $V=A\oplus B$ and
$b=b_1\oplus b_2$ along $V=A'\oplus B'$, where $a_1$ is conjugate to $b_1$ by an isomorphism
$A\to A'$ and $a_2$ to $b_2$ by an isomorphism $B\to B'$, then $a$ and $b$ are conjugate in $G$.
Products of conjugates in $\GL(A)\times\GL(B)$ can be formed in each factor separately.

\section{The cases \texorpdfstring{$n=3$}{n=3} and \texorpdfstring{$n=4$}{n=4}}\label{sec:small}

The group $\GL_3(\F)$ has order $168$ and $6$ conjugacy classes, and $\GL_4(\F)$, which is
isomorphic to the alternating group $A_8$, has order $20160$ and $14$ conjugacy classes.

\emph{The proof.} We fix one representative of each class. A certificate consists of, for each
invertible matrix, a matrix conjugating it to its representative, and, for each ordered pair
$(a,b)$ of nonidentity representatives, $n$ explicit conjugates of $a$ whose product is $b$. The
certificates were found by a computer search. Lean checks them by evaluation in its kernel; the
search itself is not part of the proof. This shows that every nonidentity element is a product of
$n$ conjugates of every nonidentity element, and Lemma~\ref{lem:more} gives $P(3)$ and $P(4)$.

\emph{The exact numbers.} For $\alpha\ne1$ let $\cn(\alpha)$ be the least $m$ such that every
nonidentity element lies in $C_\alpha^m$. By Lemma~\ref{lem:more}, every nonidentity element then
lies in $C_\alpha^k$ for all $k\ge\cn(\alpha)$. Theorem~\ref{thm:main} says $\cn(\alpha)\le n$.
Table~\ref{tab:cn} lists $\cn(\alpha)$ for every class of $\GL_3(\F)$ and $\GL_4(\F)$. It was
computed exhaustively by a separate program, working with whole classes: since $C_\alpha^k$ is a
union of classes, $C_\alpha^{k+1}$ is found by multiplying one representative of each class in
$C_\alpha^k$ by every element of $C_\alpha$. This computation is not formalized, but it agrees
with the Lean-checked bound $\cn(\alpha)\le n$.

In the table, $J_k$ is the unipotent Jordan block of size $k$, $\omega$ is the companion matrix of
$x^2+x+1$, and $C(p)$ is the companion matrix of $p$.

\begin{table}[ht]
\centering
\small
\begin{tabular}{lrrccc}
\toprule
class of $\alpha$ & order & size & J1 block & $\delta(\alpha)$ & $\cn(\alpha)$\\
\midrule
\multicolumn{6}{l}{$\GL_3(\F)$}\\
$J_2\oplus1$ (transvection) & 2 & 21 & yes & 1 & 3\\
$\omega\oplus1$ & 3 & 56 & yes & 2 & 2\\
$J_3$ & 4 & 42 & no & 2 & 2\\
$C(x^3+x+1)$ & 7 & 24 & no & 2 & 3\\
$C(x^3+x^2+1)$ & 7 & 24 & no & 2 & 3\\
\midrule
\multicolumn{6}{l}{$\GL_4(\F)$}\\
$J_2\oplus1\oplus1$ (transvection) & 2 & 105 & yes & 1 & 4\\
$J_2\oplus J_2$ & 2 & 210 & no & 2 & 3\\
$\omega\oplus1\oplus1$ & 3 & 1120 & yes & 2 & 2\\
$\omega\oplus\omega$ & 3 & 112 & no & 2 & 4\\
$J_3\oplus1$ & 4 & 1260 & yes & 2 & 2\\
$J_4$ & 4 & 2520 & no & 3 & 2\\
$C(x^4+x^3+x^2+x+1)$ & 5 & 1344 & no & 3 & 3\\
$J_2\oplus\omega$ & 6 & 3360 & no & 3 & 2\\
$C\big((x^2+x+1)^2\big)$ & 6 & 1680 & no & 3 & 2\\
$C(x^3+x+1)\oplus1$ & 7 & 2880 & yes & 3 & 2\\
$C(x^3+x^2+1)\oplus1$ & 7 & 2880 & yes & 3 & 2\\
$C(x^4+x+1)$ & 15 & 1344 & no & 3 & 2\\
$C(x^4+x^3+1)$ & 15 & 1344 & no & 3 & 2\\
\bottomrule
\end{tabular}
\medskip
\caption{The least $m$ such that every nonidentity element is a product of $m$ conjugates of
$\alpha$, for every nonidentity class of $\GL_3(\F)$ and $\GL_4(\F)$.}
\label{tab:cn}
\end{table}

Three features of the table are worth noting.
\begin{itemize}
\item The bound $n$ is attained by the transvections, as the rank argument of the introduction
predicts, and in $\GL_4(\F)$ also by $\omega\oplus\omega$. This is the fixed-point-free element of
order $3$, the same kind of matrix ($g^2+g+1=0$) that is the exceptional target in
Lemma~\ref{lem:P}.
\item Most classes need only $2$ conjugates, far below $n$.
\item In $\GL_4(\F)$ the classes $J_2\oplus J_2$ and $\omega\oplus\omega$ have no J1 block but
minimum rank $2$. So Lemma~\ref{lem:minrank} below fails for $n=4$, and its hypothesis $n\ge5$
is needed.
\end{itemize}

\section{Matrices with a J1 block}\label{sec:J1}

In this section $n\ge5$, and we assume $P(n-1)$ and $P(n-2)$.

\begin{lemma}\label{lem:J1target}
If $\alpha\ne1$ and $h\ne1$ both have a J1 block, then $h\in C_\alpha^k$ for every $k\ge n-1$.
\end{lemma}

\begin{proof}
Write $\alpha=1\oplus\alpha'$ and $h=1\oplus h'$ along splittings $\langle v\rangle\oplus U$ and
$\langle w\rangle\oplus U'$. Then $\alpha'\ne1$ and $h'\ne1$. After conjugating $\alpha$ we may
assume the two splittings coincide. By $P(n-1)$, $h'$ is a product of $k$ conjugates of $\alpha'$
in $\GL(U')$, and adding the identity on the first summand gives $h\in C_\alpha^k$.
\end{proof}

\begin{lemma}[Lemma P]\label{lem:P}
Let $\alpha\ne1$ have a J1 block and let $g\ne1$. If $\alpha^2=1$, assume also
$g^2+g+1\ne0$. Then there is $c\in C_\alpha$ such that $gc^{-1}$ has a J1 block.
\end{lemma}

\begin{proof}
\emph{Step 1.} Suppose $v,w\in V$ and $\varphi,\varphi'\in V^*$ satisfy, with
$a=\varphi(g^{-1}v)$,
\begin{equation}\label{eq:P}
\varphi(v)=1,\quad \varphi(gv)=0,\quad \varphi'(w)=a,\quad \varphi'(\alpha w)=1,\quad
\varphi'(\alpha^2w)=0.
\end{equation}
Then the vectors $g^{-1}v,\,v$ with the covectors $\varphi\circ g,\ \varphi-a\,\varphi\circ g$ form a
frame, and so do $w,\,\alpha w$ with $\varphi'\circ\alpha,\ \varphi'-a\,\varphi'\circ\alpha$. Let $t$
be given by Lemma~\ref{lem:frame}, so that $tw=g^{-1}v$, $t\alpha w=v$,
$\varphi\circ g\circ t=\varphi'\circ\alpha$ and
$(\varphi-a\,\varphi\circ g)\circ t=\varphi'-a\,\varphi'\circ\alpha$. The last two give
$\varphi\circ t=\varphi'$. Put $c=t\alpha t^{-1}$ and $h=gc^{-1}=gt\alpha^{-1}t^{-1}$. Then
\[
hv=gt\alpha^{-1}(\alpha w)=gtw=v,\qquad
\varphi(hx)=(\varphi\circ g\circ t)(\alpha^{-1}t^{-1}x)=\varphi'(t^{-1}x)=\varphi(x),
\]
so $h$ has the J1 block $(v,\varphi)$.

\emph{Step 2: $\alpha^2\ne1$.} Pick $v$ with $gv\ne v$. Then $v$ and $gv$ are linearly independent,
so there is $\varphi$ with $\varphi(v)=1$ and $\varphi(gv)=0$. Let $(u,\psi)$ be a J1 block of
$\alpha$. Then $(\alpha^2+\alpha+1)u=3u=u\ne0$, so both $\ker(\alpha^2+\alpha+1)$ and
$\ker(\alpha^2-1)$ are proper subspaces, and some $w$ lies in neither. The vectors
$w,\alpha w,\alpha^2w$ are linearly independent: none is $0$, they are pairwise distinct because
$\alpha^2w\ne w$, and their sum is not $0$. So there is $\varphi'$ with the values required
in~\eqref{eq:P}.

\emph{Step 3: $\alpha^2=1$ and $g^2+g+1\ne0$.} Since $g\ne1$, the subspaces $\ker(g-1)$ and
$\ker(g^2+g+1)$ are proper, so some $v$ lies in neither. Then $v\notin\Span(gv,g^{-1}v)$: indeed
$v\ne0$, $v\ne gv$, $v\ne g^{-1}v$, and $v\ne gv+g^{-1}v$ because applying $g$ would give
$(g^2+g+1)v=0$. Choose $\varphi$ with $\varphi(v)=1$ and $\varphi(gv)=\varphi(g^{-1}v)=0$, so
$a=0$. Pick $w$ with $\alpha w\ne w$ and $\varphi'$ with $\varphi'(w)=0$ and
$\varphi'(\alpha w)=1$. Then $\varphi'(\alpha^2w)=\varphi'(w)=0$, and~\eqref{eq:P} holds.
\end{proof}

\begin{remark}\label{rem:sharpP}
The extra hypothesis in Lemma~\ref{lem:P} cannot be dropped. Let $\alpha^2=1$ and
$g^2+g+1=0$, and suppose $c\in C_\alpha$ and $h=gc^{-1}$ has a J1 block $(v,\varphi)$. Since
$c=c^{-1}$, the equation $hv=v$ gives $cv=g^{-1}v$, and $\varphi\circ h=\varphi$ gives
$\varphi\circ g=\varphi\circ c$. Evaluating at $v$ gives $\varphi(gv)=\varphi(g^{-1}v)$. But
$g(g+1)=g^2+g=1$, so $g^{-1}=g+1$ and $\varphi(g^{-1}v)=\varphi(gv)+\varphi(v)=\varphi(gv)+1$, a
contradiction.
\end{remark}

The case excluded in Lemma~\ref{lem:P} is handled directly. In $\GL_2(\F)$, let $\sigma$ be the
swap of the two coordinates and $\omega=\left(\begin{smallmatrix}0&1\\1&1\end{smallmatrix}\right)$,
of order $3$. Then $\omega=\sigma\cdot(\omega\sigma\omega^{-1})$ is a product of two conjugates of
$\sigma$, and since $\sigma^2=1$ it is a product of $2j$ conjugates of $\sigma$ for every $j\ge1$.

\begin{lemma}\label{lem:exception}
Let $\alpha\ne1$ with $\alpha^2=1$, and let $g$ satisfy $g^2+g+1=0$. Then $n$ is even and
$g\in C_\alpha^n$.
\end{lemma}

\begin{proof}
For $v\ne0$ we have $gv\ne v$, since otherwise $v=3v=(g^2+g+1)v=0$. Choose $\psi$ with
$\psi(v)=1$ and $\psi(gv)=0$. Using $g^2=g+1$, the pairs $v,gv$ and $\psi,\psi\circ g$ form a
frame, and both $\Span(v,gv)$ and $\ker\psi\cap\ker(\psi\circ g)$ are $g$-invariant. In the basis
$v,gv$ the map $g$ acts as $\omega$, so $g=\omega\oplus g'$, where $g'$ again satisfies
$g'^2+g'+1=0$. By induction $n$ is even. In the same way, for $u$ with $\alpha u\ne u$, the plane
$\Span(u,\alpha u)$ and a complement are $\alpha$-invariant and $\alpha=\sigma\oplus\alpha'$.
After conjugating $\alpha$, both decompositions use the same splitting $V=P\oplus U$ with
$\dim U=n-2\ge3$. The map $g'$ is not the identity.

If $\alpha'\ne1$, then $g'\in C_{\alpha'}^n$ by $P(n-2)$ and $\omega\in C_\sigma^n$ because $n$ is
even, so $g=\omega\oplus g'\in C_\alpha^n$.

If $\alpha'=1$, write $g=(\omega\oplus1)(1\oplus g')$. The first factor lies in $C_\alpha^2$.
Split $U=P'\oplus U'$ with $g'=\omega\oplus g''$ as above and let $\tau=\sigma\oplus1$ on $U$.
Then $\tau\ne1$, so $P(n-2)$ gives $g'\in C_\tau^{n-2}$. The map $1_P\oplus\tau$ is conjugate to
$\alpha$, since both act as $\sigma$ on a plane and as the identity on a complement. Hence
$1\oplus g'\in C_\alpha^{n-2}$ and $g\in C_\alpha^n$.
\end{proof}

\begin{proposition}\label{prop:J1}
Let $n\ge5$ and assume $P(n-1)$ and $P(n-2)$. If $\alpha\ne1$ has a J1 block, then every
$g\ne1$ lies in $C_\alpha^n$.
\end{proposition}

\begin{proof}
If $g$ has a J1 block, use Lemma~\ref{lem:J1target}. If $\alpha^2=1$ and $g^2+g+1=0$, use
Lemma~\ref{lem:exception}. Otherwise Lemma~\ref{lem:P} gives $c\in C_\alpha$ such that $h=gc^{-1}$
has a J1 block. Then $h\ne1$, since otherwise $g=c$ would have a J1 block. By
Lemma~\ref{lem:J1target}, $h\in C_\alpha^{n-1}$, and $g=hc\in C_\alpha^n$.
\end{proof}

\section{Matrices without a J1 block}\label{sec:noJ1}

\begin{lemma}\label{lem:minrank}
Let $n\ge5$ and let $\alpha\in\GL_n(\F)$ have no J1 block. Then $\delta(\alpha)\ge3$.
\end{lemma}

\begin{proof}
Let $L$ be an algebraic closure of $\F$ and $\lambda\in L$. For $\lambda=0$ the rank of
$\alpha-\lambda$ is $n$.

Let $\lambda=1$ and suppose $\rk(\alpha-1)\le2$. The fixed space $\ker(\alpha-1)$ has dimension at
least $n-2\ge3$, so it is not contained in the image of $\alpha-1$. Take a fixed vector $v$ outside
the image and $\varphi\in V^*$ with $\varphi(v)=1$ that vanishes on the image. Then
$\varphi\circ\alpha=\varphi$, and $(v,\varphi)$ is a J1 block, a contradiction.

Let $\lambda\notin\{0,1\}$ and let $E_\mu\subseteq L^n$ be the $\mu$-eigenspace of $\alpha$. Since
the entries of $\alpha$ lie in $\F$, squaring every coordinate maps $E_\lambda$ into
$E_{\lambda^2}$. This map is injective and semilinear with respect to the Frobenius automorphism
of $L$, so $\dim E_\lambda\le\dim E_{\lambda^2}$. As $\lambda^2\ne\lambda$, the two eigenspaces
intersect in $0$, so $2\dim E_\lambda\le n$ and
$\rk(\alpha-\lambda)\ge n-\lfloor n/2\rfloor\ge3$.
\end{proof}

\begin{lemma}\label{lem:classN}
Let $n\ge3$ and let $r\in\F[x]$ be irreducible of degree $n-1$. If $u\in\GL_n(\F)$ has minimal
polynomial divisible by $r$, then its minimal and characteristic polynomials both equal
$r\cdot(x+1)$. In particular any two such matrices are similar.
\end{lemma}

\begin{proof}
The characteristic polynomial $\chi_u$ is divisible by $r$, so $\chi_u=r\,s$ with $s$ monic of
degree $1$. Since $\det u\ne0$, $s(0)\ne0$, so $s=x+1$. The minimal polynomial is $r$ or
$r(x+1)$. Since $1$ is a root of $\chi_u$, it is an eigenvalue of $u$ and a root of the minimal
polynomial, while $r$, irreducible of degree at least $2$, has no root in $\F$. So the minimal
polynomial is $r(x+1)=\chi_u$, and $u$ is similar to the companion matrix of $r(x+1)$.
\end{proof}

\begin{proposition}\label{prop:cyclic}
Assume Theorems A and B. Let $n\ge5$, let $\alpha\in\GL_n(\F)$, $\alpha\ne1$, have no J1 block, and
let $k=\lceil(n-1)/3\rceil$. Then every $g\ne1$ lies in $C_\alpha^{2k+1}$.
\end{proposition}

\begin{proof}
Write $C=C_\alpha$. By Lemma~\ref{lem:minrank}, $k$ copies of $C$ have minimum ranks summing to at
least $3k\ge n-1$, and $C$ is non-central. Theorem~A gives a cyclic matrix $D\in C^k$. Its inverse
is cyclic too: if $v,Dv,\dots,D^{n-1}v$ is a basis, so is the sequence $D^{-j}(D^{n-1}v)$,
$0\le j<n$.

Let $r$ be the minimal polynomial over $\F$ of a primitive element of $\mathbb{F}_{2^{n-1}}$. It is
monic and irreducible of degree $n-1\ge2$, and $r(0)\ne0$. The only homothety in $\GL_n(\F)$ is
$1$, so $\alpha$ and $g$ are not homotheties.

Theorem~B with $M=\alpha$ and $X$ the class of $D^{-1}$ gives $\alpha=d'x'$ with $x'$ similar to
$D^{-1}$ and $r$ dividing the minimal polynomial of $d'$. Then $x'^{-1}$ is similar to $D$, so
$x'^{-1}\in C^k$ and $d'=\alpha x'^{-1}\in C^{k+1}$.

Theorem~B with $M=g$ and $X$ the class of $D$ gives $g=dx$ with $x$ similar to $D$, so $x\in
C^k$, and $r$ dividing the minimal polynomial of $d$. By Lemma~\ref{lem:classN}, $d$ is similar to
$d'$, so $d\in C^{k+1}$ and $g=dx\in C^{2k+1}$.
\end{proof}

\section{Proof of the theorems}\label{sec:proof}

\begin{proof}[Proof of Theorem~\ref{thm:main}]
We prove $P(n)$ for all $n\ge3$ by strong induction. The cases $n=3$ and $n=4$ are
Section~\ref{sec:small}. Let $n\ge5$ and $\alpha\ne1$. If $\alpha$ has a J1 block,
Proposition~\ref{prop:J1} (which uses $P(n-1)$ and $P(n-2)$, with $n-2\ge3$) gives
$G\setminus\{1\}\subseteq C_\alpha^n$. Otherwise Proposition~\ref{prop:cyclic} gives
$G\setminus\{1\}\subseteq C_\alpha^{2k+1}$ with $k=\lceil(n-1)/3\rceil$, and $2k+1\le n$ for
$n\ge5$. In both cases Lemma~\ref{lem:more} gives $G\setminus\{1\}\subseteq C_\alpha^{k'}$ for
every $k'\ge n$.
\end{proof}

\begin{proof}[Proof of Theorem~\ref{thm:far}]
Combine Proposition~\ref{prop:cyclic} with Lemma~\ref{lem:more}.
\end{proof}

\begin{corollary}
Assume Theorems A and B. For $n\ge3$ and $\alpha\in\GL_n(\F)$, $\alpha\ne1$, the identity is a
product of $n+1$ conjugates of $\alpha$.
\end{corollary}

\begin{proof}
By Theorem~\ref{thm:main}, $\alpha^{-1}=c_1\cdots c_n$ with $c_i\in C_\alpha$, so
$1=c_1\cdots c_n\alpha$.
\end{proof}

For a single class repeated $n+1$ times, this is the $\F$ analogue of the statement that the
extended covering number is at most $n+1$, which Kn\"uppel and Nielsen prove for arbitrary
classes over fields that are not too small.

\section{The formalization}\label{sec:lean}

The Lean development accompanies this note. It consists of about 2300 lines of Lean together
with generated certificate data, and builds against a pinned version of Mathlib. Its main
theorem is
\begin{verbatim}
theorem Vaserstein.question :
    KnuppelNielsen69 -> BungerNielsen12 -> Question (ZMod 2)
\end{verbatim}
together with \texttt{Vaserstein.questionFrom}, the version with $k\ge n$ factors.

\emph{What is stated.} \texttt{Question F} is defined directly from the wording of Problem
16.21, with \texttt{SpecialLinearGroup}, its centre, and similarity by an invertible matrix. It
quantifies over all $n>2$ and all non-central $\alpha,\beta$. \texttt{KnuppelNielsen69} and
\texttt{BungerNielsen12} are the statements of Theorems A and B for $\F$, with cyclicity and
minimum rank defined as in Section 1. The minimum rank is defined through an algebraic closure,
as in the formula of Kn\"uppel and Nielsen.

\emph{What is proved.} Everything else in this note, including the cases $n=3,4$, which are
proved with no hypotheses. Their certificates, about $20{,}000$ conjugating matrices for $n=4$,
are checked with \texttt{decide +kernel}, so the check runs in Lean's kernel and not in compiled
code. The development uses no \texttt{sorry}, no \texttt{native\_decide} and no axioms beyond
Lean's standard \texttt{propext}, \texttt{Classical.choice} and \texttt{Quot.sound}.

\emph{What must be trusted.} Lean's kernel and Mathlib; the published proofs of Theorems A and
B; and the transcription of those two theorems into Lean, which was compared by hand with the
statements in the papers. An error in the transcription would make the Lean theorem a statement
about slightly different hypotheses, so this is the step a careful reader should check. The
exact numbers of Table~\ref{tab:cn} come from a separate computation and are not formalized.

\section{Discussion and open questions}\label{sec:discussion}

\begin{question}
Can Theorems A and B be formalized, so that the answer for $\F$ becomes unconditional? Theorem~B
concerns a single field and an explicit construction and seems the more accessible of the two.
\end{question}

\begin{question}
Does the same argument work over other small fields? The reduction for matrices with a J1 block
uses $\F$ in several places, for instance the test for linear independence by sums of subsets,
the identity $3u=u$ in Lemma~\ref{lem:P}, and the structure of $\GL_2(\F)$ in
Lemma~\ref{lem:exception}.
\end{question}

\begin{question}
What is $\cn(\alpha)$ for general $n$? Table~\ref{tab:cn} and Theorem~\ref{thm:far} suggest that
$\cn(\alpha)$ is much smaller than $n$ unless $\alpha$ is close to the identity or to a sum of
order-$3$ rotations. Is $\cn(\alpha)$ bounded by a constant for cyclic $\alpha$, as Lev's results
give, under conditions on the eigenvalues, for larger fields?
\end{question}

\begin{question}
Can the special cases used here be proved directly? Only the case of Theorem~A with all classes
equal is used, and Theorem~B is used for a single polynomial $q$ in each dimension.
\end{question}

\subsection*{Use of AI}
As described in Section 1, this work was produced almost entirely by an AI agent (Droid, by
Factory) under the author's direction, while testing a new computational mathematics project. The
agent carried out the computations, found the argument, searched the literature, wrote the Lean
formalization and drafted this text; the author chose the problem, set the scope and takes
responsibility for the paper. The Lean check covers every step except Theorems A and B, but it
does not cover the exposition, the literature survey or the claims about prior work, which have
not been reviewed by a specialist. Corrections are welcome.

\begin{thebibliography}{9}

\bibitem{AH}
Z.~Arad and M.~Herzog (eds.),
\emph{Products of Conjugacy Classes in Groups},
Lecture Notes in Math. \textbf{1112}, Springer, 1985.

\bibitem{BN}
F.~B\"unger and K.~Nielsen,
\emph{A matrix-decomposition theorem for $\mathrm{GL}_n(K)$},
Linear Algebra Appl. \textbf{298} (1999), 39--50,
\href{https://doi.org/10.1016/S0024-3795(99)00138-X}{doi:10.1016/S0024-3795(99)00138-X}.

\bibitem{Knuppel}
F.~Kn\"uppel,
\emph{Products of similar matrices},
Linear Algebra Appl. \textbf{233} (1996), 99--108.

\bibitem{KN}
F.~Kn\"uppel and K.~Nielsen,
\emph{The extended covering number of $\mathrm{SL}_n$ is $n+1$},
Linear Algebra Appl. \textbf{418} (2006), 634--656,
\href{https://doi.org/10.1016/j.laa.2006.03.011}{doi:10.1016/j.laa.2006.03.011}.

\bibitem{Kourovka}
E.~I.~Khukhro and V.~D.~Mazurov (eds.),
\emph{Unsolved Problems in Group Theory. The Kourovka Notebook}, No.~21, 2026,
\href{https://arxiv.org/abs/1401.0300}{arXiv:1401.0300}.

\bibitem{Lev}
A.~Lev,
\emph{Products of cyclic similarity classes in the groups $\mathrm{GL}_n(F)$},
Linear Algebra Appl. \textbf{202} (1994), 235--266.

\bibitem{Lean4}
L.~de~Moura and S.~Ullrich,
\emph{The Lean 4 theorem prover and programming language},
in: Automated Deduction, CADE 28, Lecture Notes in Comput. Sci. \textbf{12699}, Springer, 2021,
625--635.

\bibitem{Mathlib}
The mathlib Community,
\emph{The Lean mathematical library},
in: Proceedings of the 9th ACM SIGPLAN International Conference on Certified Programs and
Proofs (CPP 2020), 367--381.

\end{thebibliography}

\end{document}
